%=====================================================================
\chapter{Where to Read Further}\label{app:resources}
%=====================================================================
\normalfigures

HEC's outline for this course lists four reference books. Three of them are
cited at editions that are now between nine and nineteen years old, and one
teaches a runtime that lost security support in 2018. This appendix records
what HEC prescribed, what this course used instead, and --- in each case ---
why, so that a reviewer can see that the substitution was made deliberately
rather than by neglect.

\section{What HEC Prescribed}

\rowcolors{2}{white}{lightbg}
\begin{longtable}{@{}L{6.4cm}L{2.2cm}L{5.8cm}@{}}
\toprule
\rowcolor{primary!15}
\textbf{Reference, as HEC lists it} & \textbf{Edition} & \textbf{Status} \\
\midrule
\endhead
Laudon \& Traver, \emph{E-Commerce: Business, Technology, Society} & 13th,
2017 & Still the standard text; much newer editions exist \\
\textbf{Peacock, \emph{PHP 5 E-commerce Development}} & 2010 & \textbf{Targets
a runtime unsupported since 2018} \\
Rayport \& Jaworski, \emph{Introduction to E-Commerce} & 2007 & Predates the
mobile web \\
Schneider, \emph{Electronic Commerce} & --- & Current editions exist \\
\bottomrule
\end{longtable}

\begin{alertbox}[On the Peacock book, which is the interesting one]
It is listed as reference 2, and \textbf{the ``Course Content'' paragraph of
HEC's outline is that book's table of contents}, nearly verbatim and in order.
That finding is recorded in the Course Specification at the front of this
handout, and it explains a great deal: why the clause list reads like a build
sequence, why ``Checkout'' appears twice, and why gift vouchers are named
(\chref{discounts} found the platform has none --- they were a feature of the
2010 cart the book builds).

\smallskip
\textbf{Why this course does not follow it.} PHP 5 reached end of life on 31
December 2018 and has received no security fix since. Teaching a 2010 cart on
a dead runtime would mean teaching, as current practice, patterns that are now
vulnerabilities --- unprepared queries, \texttt{mysql\_*} functions, MD5
password hashing. \chref{accounts} measured what this platform does instead.

\smallskip
\textbf{It is still worth half an hour of your time.} Read its contents page
beside HEC's clause list, and you will understand how a curriculum outline is
made, which is a more useful lesson than anything in the chapters.
\end{alertbox}

\section{What This Course Used Instead}

\begin{examplebox}[Business and models --- \chref{what} to \chref{planning}]
\begin{tight}
  \item \textbf{Laudon \& Traver}, \emph{E-Commerce: Business, Technology,
        Society}, \textbf{current edition}. The same book HEC names, read at
        an edition that knows about mobile commerce and marketplaces. Its
        business-model taxonomy is what \chref{models} is organised around.
  \item \textbf{Pakistan Bureau of Statistics} and the \textbf{State Bank of
        Pakistan}'s payment system reviews, for market figures. Prefer a
        primary statistic you can cite to a number from a blog post.
\end{tight}
\end{examplebox}

\begin{examplebox}[The platform --- \chref{stack} to \chref{orders}]
\begin{tight}
  \item \textbf{OpenCart's own documentation} at \texttt{docs.opencart.com},
        and, more usefully, \textbf{its source}. This handout read the source
        for every claim it makes about behaviour, because the documentation is
        thin and in several places describes an older version.
  \item \textbf{The MariaDB Knowledge Base} for \texttt{EXPLAIN}, the slow
        query log and index behaviour --- \chref{catalogue}'s method comes
        from it.
  \item \textbf{The Docker documentation} on Compose, volumes and healthchecks.
        \chref{stack} and \chref{basket} rest on the volume semantics it
        defines.
  \item \textbf{The PHP manual} on \texttt{password\_hash} and sessions. It is
        unusually good and is the primary source for \chref{accounts}.
\end{tight}
\end{examplebox}

\begin{examplebox}[Security --- \chref{security}]
\begin{tight}
  \item \textbf{OWASP Top Ten} --- the list everyone cites. Read it to know the
        vocabulary, then note that \chref{security} measured the two most
        famous entries as already handled by the framework.
  \item \textbf{OWASP Application Security Verification Standard (ASVS)} ---
        far more useful than the Top Ten for actual work, because it is a
        checklist of requirements at three levels rather than a ranking of
        fashionable attacks.
  \item \textbf{OWASP Secure Headers Project}, which is where this course's
        \texttt{security.conf} comes from.
  \item \textbf{PCI DSS and the SAQ-A questionnaire}, from the PCI Security
        Standards Council. \chref{gateway} summarises it; read the actual
        questionnaire once, because it is shorter than its reputation.
  \item \textbf{Mozilla Developer Network} on Content Security Policy, which
        is the clearest writing on the subject anywhere.
\end{tight}
\end{examplebox}

\begin{examplebox}[Law and policy --- \chref{legal}]
Primary sources only. Secondary summaries of Pakistani e-commerce law are
frequently out of date and occasionally wrong about which tier of government
is responsible.

\begin{tight}
  \item \textbf{Electronic Transactions Ordinance 2002} --- what makes an
        online order enforceable.
  \item \textbf{Prevention of Electronic Crimes Act 2016} --- read sections 3
        and 4 before \chref{security}'s laboratory, not after.
  \item \textbf{E-Commerce Policy of Pakistan (Ministry of Commerce, 2019)} ---
        the policy framework, and short enough to read in one sitting.
  \item \textbf{The provincial consumer protection Acts} --- Punjab (2005),
        Sindh (2014), Khyber Pakhtunkhwa, Balochistan. Read the one your shop
        is in.
  \item \textbf{Federal Board of Revenue} for the current sales tax position,
        and the \textbf{Finance Act} of the year in question for the rate.
        \chref{shipping} deliberately refuses to be the authority for a number
        set annually.
\end{tight}
\end{examplebox}

\begin{conceptbox}[How to read any of this]
\begin{tightnum}
  \item \textbf{Prefer the primary source.} The manual, the standard, the Act.
        Every secondary account of them in this handout can be checked against
        one.
  \item \textbf{Prefer the source you can run.} \texttt{EXPLAIN} beats an
        article about indexes; the shop's own code beats its documentation.
        This is the method of the whole course.
  \item \textbf{Date everything.} A rate, a version, an end-of-life date and a
        statutory limit all expire. A claim with no date is a claim you cannot
        check.
  \item \textbf{Distrust anything about ranking, conversion or performance that
        comes without a method.} \chref{seo} declined to make claims it could
        not check, and the same standard applies to what you read.
\end{tightnum}
\end{conceptbox}

\section{A Note on Using Tutorials}

Most of what you will find searching for this platform was written for version
1.5 or 3.x and does not apply. Two findings from this course make the point
sharply:

\begin{tight}
  \item \textbf{``Rename \texttt{.htaccess.txt}''} is in every guide, and
        \chref{security} measured it as doing nothing at all on a modern Apache
        where \texttt{AllowOverride} is \texttt{None}. The advice has been
        repeated for fifteen years and has quietly stopped being true.
  \item \textbf{``Turn on SEO URLs''} is in every guide, and \chref{seo}
        measured that doing so --- without the server-side rewrite the same
        guides assume is present --- points \emph{every link in the shop} at a
        404.
\end{tight}

\smallskip
\textbf{Neither piece of advice was ever wrong. Both stopped being right
without anybody updating them.} That is the normal fate of written advice about
software, and it is why this course measures rather than cites. Read tutorials
for the idea; verify the behaviour in your own container.
